% ---
% id: receiver/rem:whyprojector
% kind: rem
% title: "Derivations without normal ordering"
% summary: >-
%   The normal-ordered form is stated for comparison and used in no proof.
%   Every operator identity from the port Grams to the emitted weight follows
%   from the functoriality of the Fock lift and vacuum elimination, block by
%   block in finite dimension, with no convergence question.
% prominence: standard
% status: accepted
% uses: [prop:normord, lem:fock-functor]
% mentions: [lem:vacelim, thm:port-grams, thm:representation-independence, thm:click-povm, thm:coarse-graining, thm:multiclick-bound, thm:emitted-bound, receiver/detector-imperfections/thm:efficiency]
% assumes: []
% refs: []
% provenance: newly-derived
% lean:
% ---

\begin{remark}
\label{rem:whyprojector}
Proposition~\ref{prop:normord} is stated because it is the form in which these
operators appear in the optics literature and in every numerical
implementation, and because it makes the efficiency substitution of
Theorem~\ref{thm:efficiency} transparent. It is not used in the proof of any claim.
Every operator identity of the sections from the port Grams to the emitted
weight is obtained from Lemma~\ref{lem:fock-functor}\,(i)--(iv) and from the
vacuum elimination of Lemma~\ref{lem:vacelim}, which involve only
finite-dimensional linear algebra on each photon-number block: no unbounded
operator, no operator-valued series and no convergence question arises
anywhere in
Theorems~\ref{thm:port-grams}, \ref{thm:click-povm}, \ref{thm:representation-independence}, \ref{thm:coarse-graining},
\ref{thm:multiclick-bound} and~\ref{thm:emitted-bound}. A submodule whose
closed forms need series or integrals, as conditioning on a probe state can,
says so where it uses them.
\end{remark}
